\chapter[Research agenda]{Open Calculations and a Research Agenda}
\label{ch:agenda}

\paragraph{Learning outcomes.}
After this chapter the reader should be able to:
\begin{enumerate}
  \item distinguish a load-bearing assumption from a critical-path
        problem and an optional refinement;
  \item pick a first calculation that actually moves the framework,
        rather than a calculation that only decorates Chapter~15;
  \item state what would have to be true before NAMCG may be
        described as finished rather than as a programme;
  \item use the appendices as tools, not as a substitute for the
        open problems listed here.
\end{enumerate}

Chapter~16 said what would kill the theory. This chapter says what
would complete it. The two lists are not the same.

\section{What is closed conditionally inside the book}
\label{sec:ch17-closed}

The following are closed \emph{as arguments from named hypotheses
to named conclusions}. They are not closed as derivations from
empty premises.

\begin{itemize}
  \item Composite metric from a non-degenerate condensate frame
        (Repair~II; Definition~\ref{def:composite-metric},
        Assumption~\ref{ass:nondeg}). The Repair~I formula of
        Definition~\ref{def:vierbein} is not a constructed frame
        (Problem~C).
  \item Algebraic split of \(\delta\Phi_{\mu\nu}\) (Chapter~5).
  \item \(m_T=0\) from entanglement equilibrium plus Fierz--Pauli
        (Theorem~\ref{thm:mT}).
  \item The conditional, truncated heat-kernel identity for specified
        sectors and scheme (Theorem~\ref{thm:Geff}); physical
        \(G_{\mathrm{eff}}\) remains open in Problem~F.
  \item Lorentz locking in unitary gauge
        (Theorem~\ref{thm:locking}).
  \item Conservation of the total stress
        (Theorem~\ref{thm:conservation}).
  \item Replacement of a bulk quartic by an infrared boundary
        datum (Theorem~\ref{thm:Leff}).
  \item A legal Newtonian laboratory for
        Assumption~\ref{ass:drag} (Chapter~15).
\end{itemize}

A reader who has worked every exercise through Chapter~16 can
reproduce that list. Reproducing it is not the same as removing
the hypotheses.

\section{Status vocabulary and completion labels}
\label{sec:ch17-labels}

Status words are \emph{OPEN} and \emph{CLOSED CONDITIONALLY}.
Closed conditionally means proved from named hypotheses to named
conclusions; it does not remove those hypotheses. Completion words
are \emph{DERIVED} and \emph{PROGRAMME-CONDITIONAL}. Derived means
that the problem's target has no remaining assumed premise in the
claimed chain. Programme-conditional means an assumption, input,
reduced scope, or unresolved branch remains.

\begin{quote}
Keeping an assumption, treating a quantity as an input, narrowing
the scope, or choosing a scope that leaves a branch unresolved
documents a limitation and keeps the programme conditional. It does
not close the problem or satisfy the finished-derivation criterion
of Section~\ref{sec:ch17-done}. Closure of a critical-path problem
means DERIVED. A reduced-scope result is stated as such and the
problem stays OPEN. A critical-path problem is met for the
matrix-derived label only by a derivation from the matrix theory; an
external justification, model input, imposed premise, or reduced
scope supports at most a scope-limited Presentation-B model.
\end{quote}

OPTIONAL means that a problem may remain open within a stated scope
provided the claim given up is stated. Optional does not mean solved.
Priority words are CRITICAL PATH and OPTIONAL.
In Chapter~16, a \emph{load-bearing assumption} is a premise used by
the chain; it is not thereby a critical-path problem.

\begin{center}
\begin{tabular}{@{}p{0.31\textwidth}p{0.63\textwidth}@{}}
\toprule
Label & Meaning \\
\midrule
\textbf{Completed infrared derivation (matrix-derived)} &
The only label for meeting the five-item criterion in
Section~\ref{sec:ch17-done}; every critical-path problem is DERIVED
for NAMCG from the matrix theory. \\
\textbf{Scope-limited Presentation-B infrared model} &
A separately stated model with its imposed premises named. It is
not matrix-derived, does not close critical-path problems, and does
not satisfy the five-item criterion. \\
\textbf{Programme} &
Any other status; one or more critical-path premises or targets
remain unresolved. \\
\bottomrule
\end{tabular}
\end{center}

\section{Critical-path problems: B, C, and D}
\label{sec:ch17-bearing}

Problems B, C, D, F, and G are critical path for a completed
matrix-derived infrared derivation. The order below follows
dependencies, not project size.

\paragraph{Problem B: entanglement equilibrium from the path
integral.}
\emph{CRITICAL PATH; OPEN.} The fixed-algebra first-law identity
is established, but not the required stationarity.
Assumption~\ref{ass:EE} is the physical engine of
Theorem~\ref{thm:mT}. Derive \(\delta S_D=0\) at fixed volume as
a Ward identity of the microscopic action~\eqref{eq:Smicro}, or
of~\eqref{eq:Smatrix} after Problem~A, without importing Jacobson's
hypothesis from a theory that already has a metric. A replica
computation on small diamonds, with the polar vacuum as background,
is the obvious route. A result that holds only for conformal
fluctuations must be flagged as such, because the tensor block is
not conformal.
Closure requires fixed-volume variations of the tensor (non-conformal)
block from the microscopic action, a specified map
\(g^{(0)}\mapsto g\), and no use of an already-existing metric. A
conformal-only derivation is a sub-case and leaves B OPEN.

\paragraph{Problem C: the conformal mode.}
\emph{CRITICAL PATH; OPEN.} The radial potential curvature is
known, but the physical scalar operator and coupling are not.
Assumption~\ref{ass:gap-psi} is the remaining scalar-tensor
ceiling. Compute the full quadratic scalar operator on the locked
vacuum in a constructed frame: the canonically normalised kinetic
coefficient of the pure-trace scalar (including the induced
Einstein--Hilbert contribution), its mixing with \(T_{\mu\nu}\), its
non-minimal curvature coupling, and its matter coupling. The
potential~\eqref{eq:V} alone fixes only the radial curvature
\(\lambda v^{2}\). From the operator, produce the pole mass, or a
ratio \(m_\psi/v\), together with the coupling that can be compared
with Cassini bounds. If \(m_\psi\) is not parametrically of order
\(\sqrt{\lambda}\,v\), the assumption is false and the infrared
theory must be rewritten as scalar-tensor NAMCG.
Closure requires the operator, scalar coupling, and phenomenological
comparison that define and test either the safe-mass or scalar-tensor
outcome. Declaring a scalar-tensor scope without those data leaves
Problem~C open and the result programme-conditional.

\paragraph{Lemma (tree-level radial mass).}
\label{lem:msigma}
On the locked vacuum \(\langle\Phi^{I}{}_{A}\rangle=v\delta^{I}{}_{A}\),
write the polar field
\(\Phi=(v+\sigma)\Lambda\) with \(\Lambda^{\dagger}\Lambda=\mathbb{I}\)
and \(\langle\sigma\rangle=0\).
Then \(\Phi^{\dagger}\Phi=(v+\sigma)^{2}\mathbb{I}\), and
\eqref{eq:V} collapses to a singlet potential.
With \(N=\mathrm{Tr}(\mathbb{I})\) locked directions,
\begin{equation}
  V
  =\frac{\lambda}{4}N\bigl[(v+\sigma)^{2}-v^{2}\bigr]^{2}
  =N\lambda v^{2}\,\sigma^{2}+O(\sigma^{3}).
\end{equation}
The kinetic term in~\eqref{eq:Smicro}, in the signature \((-,+,+,+)\),
contributes \(-N\,(\partial\sigma)^{2}\).
Canonical normalisation \(\sigma_{c}=\sigma\sqrt{2N}\)
puts the quadratic action in the form
\(-\tfrac12(\partial\sigma_{c})^{2}-\tfrac12 m_{\sigma}^{2}\sigma_{c}^{2}\),
with
\begin{equation}
  \label{eq:msigma}
  m_{\sigma}^{2}=\lambda v^{2},
  \qquad
  m_{\sigma}=v\sqrt{\lambda}.
\end{equation}
The would-be Goldstones in \(\Lambda\) remain massless at this order.
They are the modes Chapter~10 is required to eat.
They are not \(\psi\).

\paragraph{Lemma (radial potential curvature in the pure-trace sector).}
\label{lem:psi-sigma}
Work in the pure-trace sector: \(T_{\mu\nu}=0\), \(F_{\mu\nu}=0\),
and \(\omega_{\mu\nu}=0\). Two repairs of the frame must be
kept distinct.

\emph{Repair II, unitary gauge.}
Chapter~10 permits \(\Lambda=\mathbb{I}\) wherever
Assumption~\ref{ass:nondeg} holds. The vacuum is then the frame,
\(e^{A}{}_{\mu}=\Phi^{A}{}_{\mu}/v\), so
\begin{equation}
  \label{eq:repairII-trace}
  g_{\mu\nu}
  =(1+\sigma/v)^{2}\eta_{\mu\nu},
  \qquad
  \psi
  \equiv\eta^{\mu\nu}h_{\mu\nu}
  =8\sigma/v+O(\sigma^{2}).
\end{equation}
The conformal mode and the radial singlet are the same field in
this gauge. The potential~\eqref{eq:V} gives the radial singlet the
curvature \(V''\propto\lambda v^{2}\), which with the microscopic
kinetic normalisation of Lemma~\ref{lem:msigma} is
\(m_\sigma^{2}=\lambda v^{2}\). This is a statement about the potential
only. It is not yet the mass of a canonically normalised pole of
\(\psi\) (see the identification paragraph below).

\emph{Repair I.}
Definition~\ref{def:vierbein} does not define a vierbein
(Appendix~\ref{app:status}, Section~\ref{sec:status-C}), so \(\psi\)
has no definition in terms of \(\Lambda\) and no Hessian on
\((\sigma,\psi)\) is computed for Repair~I. What survives is that
\eqref{eq:V} depends only on
\(\Phi^{\dagger}\Phi=(v+\sigma)^{2}\mathbb{I}\), so it gives
curvature only to \(\sigma\) and none to the
\(\mathfrak{so}(3,1)\)-valued directions of \(\Lambda\).

\emph{Identification \(g^{(0)}=g[e]\).}
For the metric~\eqref{eq:repairII-trace}, \(g=\Omega^{2}\eta\) with
\(\Omega=1+\sigma/v\), the induced Einstein--Hilbert density is
\(\sqrt{-g}\,R=-6\,\Omega\,\Box\Omega\), which equals
\(6(\partial\Omega)^{2}\) up to a total derivative. With the curvature
convention of Chapter~2, Section~\ref{sec:ch2-manifolds} (identified in
Chapter~9 as the one for which \(R>0\) on spheres), this adds to the
Lagrangian \(+\,3(\partial\sigma)^2/(8\pi G_{\mathrm{eff}}v^{2})\)
if the Einstein--Hilbert term has the form of~\eqref{eq:W-EH} with
\(G_{\mathrm{eff}}>0\); that sign of \(G_{\mathrm{eff}}\) is a declared
condition (Chapter~8, Problem~F), not a derived result. In the signature
\((-,+,+,+)\) a healthy scalar enters as
\(-\tfrac12 Z(\partial\sigma)^2\) with \(Z>0\), so this term has the
conformal-factor (ghost) sign relative to the microscopic term
\(-N(\partial\sigma)^2\) of Lemma~\ref{lem:msigma}. Conditional on the
pure-trace truncation, on adding the induced term to the microscopic one,
and on \(G_{\mathrm{eff}}>0\) in the Lorentzian continuation, the quadratic
action
\(-\tfrac12 Z(\partial\sigma)^2-\tfrac12 V''\sigma^2\) has
\begin{equation}
  \label{eq:Z-sigma}
  Z=2N-\frac{3}{4\pi G_{\mathrm{eff}}v^{2}}
   =2N\Bigl(1-\frac{3M_{\mathrm P}^{2}}{Nv^{2}}\Bigr),
  \qquad
  V''=2N\lambda v^{2},
  \qquad
  M_{\mathrm P}^{2}=\frac1{8\pi G_{\mathrm{eff}}}.
\end{equation}
If \(Nv^{2}>3M_{\mathrm P}^{2}\) then \(Z>0\) and a pole of squared mass
\(m^{2}=V''/Z=\lambda v^{2}\bigl(1-3M_{\mathrm P}^{2}/(Nv^{2})\bigr)^{-1}\)
exists in this truncation; it tends to~\eqref{eq:msigma} only when
\(3M_{\mathrm P}^{2}/(Nv^{2})\to0\), i.e.\
\(G_{\mathrm{eff}}Nv^{2}\to\infty\). For \(Z<0\) (i.e.\
\(Nv^{2}<3M_{\mathrm P}^{2}\)) the truncated mode has a wrong-sign kinetic
term and, with \(\lambda>0\), also \(m^{2}=V''/Z<0\): in this truncation it
is both ghost-like and tachyonic, so the kinetic sign is not the only
instability. The pure-trace truncation therefore cannot support
Assumption~\ref{ass:gap-psi} there; this signals an instability of the
chosen background and truncation and is not a statement about the full
scalar spectrum. At \(Nv^{2}=3M_{\mathrm P}^{2}\) one has \(Z=0\): the
quadratic kinetic term of this truncation degenerates, and the formula
establishes no canonically normalized finite-mass pole there. Since
\(v/M_{\mathrm P}\) requires \(G_{\mathrm{eff}}\), this ties Problem~C to
Problem~F. The mixing with \(T_{\mu\nu}\), the non-minimal curvature
coupling, and the matter coupling are not derived here, and no Cassini
conclusion follows.

\paragraph{Cassini, revised.}
The potential~\eqref{eq:V} gives the radial singlet the curvature
\(\lambda v^{2}\) in the Repair~II pure-trace sector and produces no
massless scalar by itself. The bound
\(\gamma-1=(2.1\pm 2.3)\times 10^{-5}\)
\cite{Bertotti2003}
constrains a massless or long-range scalar partner. Whether the
physical scalar of this sector is short-range on an astronomical
unit depends on the total kinetic coefficient, the non-minimal
coupling and the matter coupling, none of which is computed. This is
not a Cassini bound on NAMCG.

\paragraph{Status of Assumption~\ref{ass:gap-psi}.}
Not promoted. What is established is the potential curvature
\(m_\sigma^{2}=\lambda v^{2}\) (Lemma~\ref{lem:msigma}) and its
identification with the conformal mode in the Repair~II pure-trace
sector, under Assumptions~\ref{ass:frame-lock} and~\ref{ass:nondeg}.
The physical scalar spectrum, the mixing with \(T_{\mu\nu}\) and the
matter coupling remain open (Appendix~\ref{app:status},
Section~\ref{sec:status-C}). Loop corrections and the wrong-sign
conformal kinetic term of the heat kernel remain Chapter~8
bookkeeping.


\paragraph{Problem D: local Lorentz invariance from Presentation~A.}
\emph{CRITICAL PATH; OPEN.} Generator counting is closed
conditionally, while local Lorentz symmetry is imposed.
Assumption~\ref{ass:local-lorentz} puts
\(\mathrm{SO}(3,1)\) into Presentation~B by hand. Either derive a
local Lorentz gauge symmetry from the matrix algebra and condensate
(D1), or state that NAMCG is a theory of a condensate field on an
auxiliary spacetime and stop calling Presentation~A ontologically
prior (D2). D1 closure requires the \(\omega\), \(A_\mu\), \(B\), and
ghost operators and the resulting field content to be specified.
Under D2, even with these operators specified, the result is only a
scope-limited Presentation-B model: D remains OPEN and CRITICAL PATH,
and a full matrix-derived spectrum is not established. The Chapter~10
generator count identifies six
Goldstone orbit variables removable in unitary gauge; it does not
determine the dynamics of the independent diagonal connection \(B\)
or whether torsion vanishes, propagates, or is massive. Torsion-free
\(B\) follows only under the conditional Palatini--Holst hypotheses
of Proposition~\ref{prop:B-torsion-free}; NAMCG has not been shown
to induce them. For a Laplace-type operator,
\(a_1=\operatorname{tr}(R(g)/6-E)\); its explicit universal
curvature term is the Levi-Civita scalar \(R(g)\), while unspecified
\(E\) leaves additional \(B\)-dependence open and does not establish
a Palatini \(R(B)\) term. If \(B\) enters the actual bundle
connection, \(a_2\) may contain \(\operatorname{tr}F(B)^2\); its
survival in the weighted trace depends on unspecified operators and
weights. Deriving the connection and endomorphism dependence or
counting \(B\) and torsion as additional field content remains open;
Problem~D is unresolved.
Within D2, reduced claims may be limited to Levi-Civita metric
geometry of \(g=e^T e\) and conservation of the total stress, without
using \(B\); unitary-gauge Lorentz locking removes the six Goldstone
orbit variables, while frame/metric bookkeeping remains provisional.
D2 does not establish a full-spectrum account, the dynamics or
masslessness of \(B\), torsion-free \(B\), or the quadratic operator
of \(A_\mu\). Problem~F then applies only to an explicitly stated
frame/metric truncation omitting \(B\) and \(A_\mu\), conditional on
that truncation and with the corresponding ghost sector specified.

\section{Open problems by scope and priority}
\label{sec:ch17-coeff}

\subsection{Optional problems A and E}

\paragraph{Problem A: the infrared dictionary.}
\emph{OPEN; OPTIONAL only if Presentation~B is the operative theory
and no claim is made that NAMCG follows from Presentation~A or that
Presentation~A is ontologically prior.} The claim given up is a
derivation from the matrix model; Chapter~8 remains a Presentation-B
calculation. A claim of matrix-model priority makes A critical path.
Assumption~\ref{ass:dictionary} is the map
\begin{equation}
  \label{eq:dict-again}
  [\hat X_\mu,\Phi]
  \;\longleftrightarrow\;
  i\Theta_{\mu\nu}D^\nu\Phi
\end{equation}
plus higher-derivative corrections. A derivation must specify the
Hilbert space, Weyl map, continuum limit, and error at each order in
\(\Theta\).

\paragraph{Problem E: Lorentzian signature.}
\emph{OPEN; OPTIONAL within a scope that retains
Assumption~\ref{ass:lorentzian}.} The claim given up is that
Lorentzian signature is emergent or derived. Euclidean positivity
and a formal Wick rotation do not select it; rotating because the
infrared world is Lorentzian is not a derivation.

\subsection{Critical-path items F and G}

\paragraph{Problem F: the rationals \(A,B,C\).}
\emph{CRITICAL PATH; OPEN.} The conditional heat-kernel identity
does not compute the physical coupling.
The target is to evaluate the weighted \(a_1\) and the
curvature-linear part of \(a_2\) on the locked block, including
diffeomorphism and local-Lorentz ghosts and the radial-mode
threshold, with a scheme stated. The finite-dimensional fibre count
alone does not determine these coefficients: the quadratic operator,
its connection and endomorphism, the determinant weights, and finite
counterterms enter.

\paragraph{Verified heat-kernel bookkeeping (not closure).}
For the proper-time convention of Appendix~\ref{app:seeley-dewitt},
\[
W_E=-\frac12\int_{\LUV^{-2}}^{\mu^{-2}}\frac{\mathrm{d}s}{s}\,
\mathrm{Tr}\,e^{-s\mathcal O},
\]
termwise integration of the truncated, unresummed small-\(s\)
expansion gives the \(a_0,a_1,a_2\) factors
\(-1/(64\pi^2),-1/(32\pi^2),-1/(32\pi^2)\), respectively, with
powers \(\LUV^4-\mu^4\), \(\LUV^2-\mu^2\), and
\(\ln(\LUV^2/\mu^2)\). These are not exact finite-cutoff
integrations of massive determinants; resumming mass exponentials
produces sector-dependent threshold functions. The standard
curvature-square part of \(a_2\) for \(\mathcal O=-D^2+E\) is
\[
\frac1{180}\left(R_{\mu\nu\rho\sigma}R^{\mu\nu\rho\sigma}
-R_{\mu\nu}R^{\mu\nu}+\frac52R^2\right),
\]
not the previously displayed \(7R^2/2\).

For one minimally coupled real bosonic Laplace-type field with
\(E_R=0\), the universal \(a_1\) contribution to the Euclidean
coefficient of \(\int\sqrt{g}\,R_E\) is
\[
C_{R,E}^{(a_1)}
=-\frac{N_0}{192\pi^2}(\LUV^2-\mu^2),
\]
where \(N_0\) is the fibre trace of this single sector and the
displayed limits are assigned to it. For a specified collection,
the corresponding Euclidean contribution is the sum over the
explicitly weighted sectors and their assigned limits. This is not
the full locked-block result and does not establish a ghost-subtracted
trace. It makes no claim about the Lorentzian Einstein--Hilbert sign
or a physical \(G_{\mathrm{eff}}\); those require the specified
sector sum and an explicit action/curvature continuation convention.

If \(E=E_0+QR+\cdots\), the \(a_1\) curvature coefficient is
\(\operatorname{tr}(\tfrac16\mathbb I-Q)\). For a scalar with
\(E=m^2+\xi R\), the mass-dependent curvature-linear term in \(a_2\)
is \((\xi-\tfrac16)m^2R\), vanishing at \(\xi=1/6\). Thus the
conformal radial scalar alone does not establish a nonzero logarithmic
\(B\); the additional frame operator must be specified and evaluated.

\paragraph{Missing step.}
Specify the gauge-fixed quadratic Hessian for every locked-block
field, including bundle connections, endomorphisms, diffeomorphism
and local-Lorentz ghost operators and weights, the radial threshold
prescription, finite counterterms, and continuation convention. Then
compute the weighted \(a_1\) and relevant \(a_2\) terms in that scheme
and match them to \(A,B,C\). The present manuscript does not specify
those data. The asserted net \(+2\) from the on-shell TT count is not
an off-shell heat-kernel trace, and no replacement ghost count can be
inferred without the actual gauge fixing and ghost operators. The
values \(A=1/\pi\), \(B=8\pi/3\), and \(C=0\) are not established.
Problem~F remains open.


\paragraph{Problem G: infrared boundary entropy.}
\emph{CRITICAL PATH; mechanism CLOSED CONDITIONALLY, derivation OPEN.}
Evaluate \(S_{\mathrm{ent}}^{\mathrm{(IR)}}\) on an explicit
cosmological or large-tube boundary and produce a candidate
\(\Lambda_{\mathrm{eff}}\) from~\eqref{eq:Leff}. A symbolic
formula in terms of a late-time entanglement scale is more useful
than a claim to have obtained \(10^{-122}\).

\paragraph{G completion levels.}
For G-model, a scaling or value computed within a stated assumed
boundary model is PROGRAMME-CONDITIONAL and does not satisfy the
matrix-derived completion criterion. For G-external, independently
justified boundary entropy and stationarity may support only a
scope-limited Presentation-B model; G remains open for the
matrix-derived label. G is DERIVED only when the entropy density and
stationarity relation are derived from the matrix theory and
\(\Lambda_{\mathrm{eff}}\) follows from them. A scaling selected to
reproduce a target remains an input. The required entropy density is
the one in Assumption~\ref{ass:Sent-density}, derived by the replica
construction from the matrix theory, and boundary stationarity
Assumption~\ref{ass:bdy-EE} must likewise be derived rather than
retained.

\subsection{Optional problems continued}

\paragraph{Problem H: interpolating function.}
\emph{OPEN; OPTIONAL research within the theory.} The claim given
up is a derivation of \(\nu_{\sqrt{}}\) from modular data; current
empirical comparisons remain coefficient checks.
Derive \(\nu_{\sqrt{}}\) in~\eqref{eq:nu-sqrt}, or a function
that beats it on a SPARC hold-out, from the modular data of
Chapters~6 and~12. The target is
\begin{equation}
  \label{eq:nu-target}
  \nu_{\sqrt{}}(g_{\mathrm{N}})
  =\frac{1}{1-\exp\bigl(-\sqrt{g_{\mathrm{N}}/g_s}\bigr)},
\end{equation}
which is identically \(1+n_{\mathrm{BE}}(\sqrt{g_{\mathrm{N}}/g_s})\).
Appendix~\ref{app:empirical} already shows that this law, frozen
at the published \(g_\dagger\), is equivalent to the RAR and
about \(0.2\%\) behind simple MOND, while baryons-only fails.
That ranking is a coefficient check. It does not solve
Problem~H.

The first finite computation is the replica free energy of a
free field on a weak-field diamond whose infrared boundary is
specified as part of the calculation. Chapter~\ref{ch:hbev} supplies
\(\Lambda_{\mathrm{eff}}\) but does not define a stationary boundary
geometry or its surface gravity \(\kappa_{\mathrm{IR}}\); a
placeholder \(\kappa_{\mathrm{IR}}\) with \(\Lambda=\Lambda_{\mathrm{eff}}\)
must be constructed and normalized there. Read the source-dependent
modular excitation gap \(\Delta k(g_{\mathrm{N}},\kappa_{\mathrm{IR}})\)
when a Newtonian source of acceleration \(g_{\mathrm{N}}\) is introduced;
the absolute lowest eigenvalue is not used, because it depends on an
additive normalization (Section~\ref{sec:status-H}). Problem~H is solved
if the response to the extra stress, computed from the weighted
multimode spectral sum of Section~\ref{sec:status-H} (or from a
justified single-ladder response), reproduces~\eqref{eq:nu-target}
with no SPARC input. In the single-ladder case this requires
\begin{equation}
  \label{eq:gap-H}
  \Delta k
  =
  \sqrt{\frac{g_{\mathrm{N}}}{g_s}},
  \qquad
  g_s=c_\star\,\kappa_{\mathrm{IR}},
\end{equation}
with the order-one constant \(c_\star\) obtained from the calculation
and not from galaxy data. The values \(c_\star=1\) and
\(c_\star=1/(2\pi)\) both appear in this manuscript; neither is
adopted, and a scaling
\(\Delta k\propto\sqrt{g_{\mathrm{N}}/\kappa_{\mathrm{IR}}}\) alone does
not fix \(c_\star\). A computed spectrum that does not reproduce
\eqref{eq:nu-target} does not solve Problem~H. A linear gap
\(\Delta k\propto g_{\mathrm{N}}/\kappa_{\mathrm{IR}}\)
reproduces the retired \(\nu_{\mathrm{exp}}\) and kills
Assumption~\ref{ass:sqrt-mod}.

Alternatively, prove that no local
\(\mu_{\mathrm{ent}}(g_{\mathrm{N}})\) exists and write the
nonlocal kernel instead. Refitting \(\alpha\), \(r_s\), or
\(g_s\) to rotmod files is not a solution. Identifying
\(g_s\) with \(g_\dagger\) inside the replica calculation is
not allowed; compare \(\kappa_{\mathrm{IR}}\) to \(g_\dagger\)
only after the source-dependent gap and response are calculated.

\paragraph{Problem I: dynamical origin of the drag.}
\emph{OPEN; OPTIONAL falsifiable phenomenological extension.}
It becomes critical only if the drag is claimed to be derived or
predicted. Otherwise the microscopic exponent and rate remain open,
and Assumption~\ref{ass:drag} is a working law.
Replace Assumption~\ref{ass:drag} by a low-frequency response
calculation of \(\Gamma_{ij}(\mathbf r;\ell_A,\ell_B)\) from
specified microscopic dynamics and a specified stationary reference
state, including the finite-size factors \(\ell_A,\ell_B\) and the
longitudinal and transverse components. Report the exponent and the
coefficient. Then, and only then, estimate \(\eta\) at
tidal radii and at LIGO-band separations. The static tail and its
assumptions are recorded in Problem~\ref{sec:status-I}; no static-tail
power determines the retarded response.

\paragraph{Problem J: spherical ECG solutions.}
\emph{OPEN; OPTIONAL research within the stated model scope.} The
claim given up is singularity resolution derived from the matrix
theory; reduced results remain conditional on symmetric reduction
and the ECG model assumption.
For the explicit invariant~\(\mathcal P\) in~\eqref{eq:ECG}, used in
the action~\eqref{eq:S-ECG}, Section~\ref{sec:status-J} records
checked reduced Euler--Lagrange
identities, the resulting first integral including
\(\Lambda_{\mathrm{eff}}\), its conditional regular-centre constraint
for every \(\Lambda_{\mathrm{eff}}\), the maximally symmetric
curvature relation, and the large-radius tail. These results are
conditional on the static spherical reduction capturing the full
field equations; that symmetric-criticality step is not established
here. The Einstein-like linearized spectrum on maximally symmetric
backgrounds is cited separately and does not establish nonlinear or
sourced stability. Problem~J remains open: justify the reduction
against the full field equations, establish the cubic truncation and
its microscopic matching, and construct a regular massive core from
specified sourced or modified equations, including its entanglement
stress and residual relative to the phenomenological
profile~\eqref{eq:f-reg}.

\section{Problems that look like work and are not}
\label{sec:ch17-not}

The following activities are useful as exercises and do not move
the framework.

\begin{itemize}
  \item Producing further Newtonian integrators with new colour
        maps of Figure~\ref{fig:lab}.
  \item Fitting SPARC with an unconstrained \(\mu_{\mathrm{ent}}\),
        a per-galaxy \((\alpha,r_s)\), or a new \(\nu\) that is
        not~\eqref{eq:nu-target} and is not derived from
        \(I(A:B)\). Appendix~\ref{app:empirical} already did the
        legitimate coefficient ranking.
  \item Writing \(\mathbf{a}_{\mathrm{irrev}}\) with additional
        free exponents and calling each a prediction.
  \item Replacing \(G_{\mathrm{toy}}\) by \(G_{\mathrm{N}}\) in
        a slide.
  \item Adding adjectives (``holographic,'' ``non-Abelian,''
        ``entropic'') to an abstract without adding a lemma.
\end{itemize}

Chapter~15 exists so that a legal laboratory is on the record.
It does not exist so that the laboratory can become the theory.

\section{A recommended order of work}
\label{sec:ch17-order}

The critical-path dependency order is:

\begin{enumerate}
  \item Problem~D1: derive local \(SO(3,1)\) from the matrix algebra
        and condensate; specify the \(\omega,A_\mu,B\) and ghost
        operators needed by Problem~F.
  \item Problem~C: compute the tree-level scalar operator, including
        kinetic coefficient, mixing, mass and testable matter coupling.
  \item Problem~F: compute \(A,B,C\) with the specified operators,
        ghosts, thresholds and scheme, and derive \(G_{\mathrm{eff}}\).
  \item Problem~B: derive the fixed-volume tensor-block stationarity
        from the microscopic action, not only for conformal variations.
  \item Problem~G: derive \(S_{\mathrm{ent}}^{\mathrm{(IR)}}\) and
        boundary stationarity from the matrix theory, then derive the
        \(\Lambda_{\mathrm{eff}}\) scaling or value.
\end{enumerate}

\paragraph{Alternate scope: model work outside the critical-path sequence.}
Taking Presentation~B with imposed local Lorentz symmetry (D2), an
externally justified boundary result (G-external), or \(G_{\mathrm{eff}}\)
as an input may make a scope-limited Presentation-B model calculable
and testable, for example by specifying full quadratic and ghost
operators and computing Problem~F for a stated truncation. Such work
leaves Problems~D, F, G and the matrix-derived completion criterion
open, must name its imposed premises, and is not a step in the
sequence above.

Problems~A, E, H, I, and J are optional within their stated scopes.
Appendix~\ref{app:empirical} records the coefficient-level check for
Problem~H; further galactic work should not be described as a
derivation from \(I(A:B)\) without the modular calculation.

\section{What ``finished'' would mean}
\label{sec:ch17-done}

The label \emph{Completed infrared derivation (matrix-derived)}
applies only if all five critical-path problems are DERIVED:

\begin{enumerate}
  \item Problem~B is a lemma: Assumption~\ref{ass:EE} is no longer an
        assumption, and fixed-volume variations include the
        non-conformal tensor block.
  \item Problem~C has a computed scalar operator, coupling, and
        phenomenological comparison defining and testing either a
        sufficiently heavy \(m_\psi\) or the scalar-tensor outcome.
  \item Problem~D1: an \(SO(3,1)\) gauge symmetry is derived from the
        matrix algebra and condensate, with the \(\omega,A_\mu,B\)
        and ghost operators specified and the field content stated.
  \item Problem~F has \(A,B,C\) and the sign/scheme dependence of
        \(G_{\mathrm{eff}}\) computed from the specified operators;
        \(G_{\mathrm{eff}}\) is not an input.
  \item Problem~G-derived: the infrared boundary entropy density
        and stationarity relation are derived from the matrix theory,
        and the \(\Lambda_{\mathrm{eff}}\) scaling or value follows
        from them. G-model and G-external do not meet this item.
\end{enumerate}

If local Lorentz symmetry is instead imposed (D2), even with the
full quadratic and ghost content specified, the result is a
scope-limited Presentation-B infrared model. It may be studied and
tested as such, but Problem~D stays OPEN and CRITICAL PATH and the
five-item criterion is not met. No label of ``completed infrared
derivation'' is allowed for that scope. Any critical-path item met
only by an imposed premise, input, or reduced scope likewise yields
at most a scope-limited model or programme, never the matrix-derived
label.

Problems~A, E, H, I, and J may remain open in a completed derivation
within the stated scopes; the claims given up are listed above.
Problem~A is additionally required if the theory is claimed to
follow from Presentation~A. Problem~I is a prerequisite only if
drag is claimed to be derived or predicted.

Until all five items are derived, the honest public description is
the subtitle of the book, read as a programme:
\emph{an attempt} to derive general relativity and massless
spin-2 dynamics from a matrix condensate held at entanglement
equilibrium. The attempt is written so that it can fail in
public. That is the point of Chapter~16.

\section{Exercises}
\label{sec:ch17-exercises}

\begin{enumerate}
  \item \textbf{Status report.}
    Copy the eleven-line chain of Section~\ref{sec:ch16-chain}
    and mark each line ``lemma,'' ``theorem from assumptions,''
    or ``assumption.'' Do not use a mark you cannot defend
    from the text.
  \item \textbf{Proposal.}
    Write a two-page project proposal for Problem~C or Problem~F
    only. Include the first calculation you would perform and
    the result that would make you stop.
  \item \textbf{Not a proposal.}
    Take any project that Section~\ref{sec:ch17-not} would
    reject and explain, in six sentences, why it does not
    change the status of Theorem~\ref{thm:mT}.
  \item \textbf{Public sentence.}
    Write one sentence suitable for an Amazon description that
    survives Section~\ref{sec:ch16-not-made} and the last
    paragraph of Section~\ref{sec:ch17-done}.
\end{enumerate}

\section{Bridge to the appendices}
\label{sec:ch17-bridge}

The main text is complete. Appendix~A records the Seeley--DeWitt
algorithm at the working thickness used in Chapter~8.
Appendix~B records the replica formula and the modular Hamiltonian
at the thickness used in Chapter~6. Appendix~C records the
Verlet laboratory and the energy identity used in Chapter~15.
Appendix~D records the SPARC ranking of \(\nu_{\sqrt{}}\) against
the RAR, simple MOND, and baryons-only. None of the appendices
introduces a new physical hypothesis. If an appendix appears to
do so, the hypothesis belongs in the main text or it does not
belong in the book.